\section*{Cronograma}
\justifying
Siempre que se prepara un proyecto, resulta muy importante agregar un cronograma
y un diagrama para conocer los periodos de desarrollo para el proyecto explicando
detalladamente las etapas del proyecto. En este desarrollo, deben presentarse,
agrupadas en bloques, las actividades que el alumno y su grupo de trabajo
desarrollarán. Este cronograma debe acompañarse por una sección breve a manera
de introducción y explicando su contenido, o bien hacer referencia al cuadro
mostrado en la sección de metodología.

Lo más conveniente es presentar los grupos de actividades por periodos utilizando
unidades de tiempo similares como las semanas o los meses. Así, los lectores
pueden entender rápidamente el tiempo total que abarcará el proyecto y
cuanto tiempo llevará cada etapa del mismo.
